\section{耶利米- 「耶和华建立」或「差派」}
http://www.wccec.org/newWCCEC/Version1/gym/marathon/Jeremiah.pdf
\subsection{信息二十篇}
\subsubsection{被擄前的信息二十篇}
\begin{table}[H]
\centering
\begin{tabular}{p{2.cm}|p{15.50cm}}\hline\hline
\multicolumn{1}{c|}{经文}&\multicolumn{1}{c}{内容}\\\hline
(2-3:5)&神兒女的二件惡事：離棄神活水的泉源，為自己鑿池子，卻是破裂不能存水的漏池。主是葡萄樹，我們是枝子，支取主的生命汁漿長大豐盛。神一再呼召祂的兒女歸向祂，親近祂的人，祂也必親近他，且能順服祂，抵擋魔鬼(雅四8)。
\\\hline
(3:6-6章)&神對背道兒女愛的呼喚，叫他們承認罪孽，應許有一救主要牧養他們，雖然仇敵兇殘，卻不能越過神的範圍，撒但是神手中的篩子，藉此潔淨祂的兒女。\\\hline 
 (7-10章)&脫離虛假的聚會，改正行動作為，對付罪惡，在心靈與真實裡敬拜神，就是為主而活。\\\hline
 (11-12章)&神兒女毀約帶來咒詛，回憶以往神帶領他們出埃及，脫離魔鬼的權勢，現今竟不聽從神的約，隨從自己頑梗的噁心去行，惹動神的怒氣，自作自受，雖呼求神也徒然。\\\hline
 (13章)&離棄神而獨立，遭被擄，靈裡昏盲。\\\hline
 (14-15章)&離棄神，生命枯死，失去神的同在，被擄的命運已經定奪。\\\hline
 (16-17:18)&被擄的悲哀，呼求神作避難所。\\\hline
 (17:19-27)&呼籲安息在神的手中，向己死，向神活，分別為聖單單屬於神。\\\hline
 (18章)&信而順服的呼召，敗壞背離神的國，建立順服神的國。\\\hline
 (19-20章)&裡外之火焚燒先知，又焚燒神的兒女，頑梗到底還是不悔改，將神當作平常，結局悲慘。\\\hline 
 (21章)&生命與死亡之路的選擇，順服神手中的懲罰工具巴比倫王存活？還是逃亡而死？\\\hline
 (22章)&數算猶大王室之罪，違背律法的都是罪。\\\hline
 (23:1-8)&趕散殘害羊群的牧人之罪，為己權益不惜犧牲羊群，分門別類，遭神棄絕。\\\hline
 (23:9-40)&先知為己利益妄自發預言，引導百姓走入岐路，遭神反對。\\\hline
 (24章)&呼籲民一心歸向神，如同極好的無花果，有認識神的心，一心歸向神，順服神服在巴比倫王的權勢下受管教。認識神的秘訣是一心歸向神(林後三16)。\\\hline 
(25章)&呼籲離開惡道，與神同住，被擄七十年的命運已定。\\\hline
 (26章)&再次呼籲各人回頭，離開惡道，改正行動作為挽回神的怒氣。\\\hline
(27-28章)&呼籲順服神手中的篩子巴比倫王。\\\hline
(29章)&先知寄信去被擄之地，勉勵他們順服七十年，末後必有指望，專心尋求祂的就必尋見。\\\hline
(30-33章)&管教的目的是為了拯救，拆毀是為栽植，祂所疼愛的祂必管教，被主管教的人有福了，主耶穌的死而復活必買贖祂的教會，如同先知用十七舍克勒銀子買回了地。十七是復活之數，主在17日清晨復活。\\\hline
\end{tabular}
\end{table}

\subsubsection{被擄中的信息和先知在埃及的信息}
\begin{table}[H]
\centering
\begin{tabular}{p{2.cm}|p{12.50cm}}\hline\hline
\multicolumn{1}{c|}{经文}&\multicolumn{1}{c}{内容}\\\hline
39-42章&傳達神勸告百姓投降巴比倫王，若不聽從逃去埃及必死亡。\\\hline
 (43-52章)&先知被擄埃及，指明埃及地的逃亡猶太人必亡於巴比倫。外邦十國必遭審判。\\\hline
 (46章)&主前605年，迦基米施之役，埃及敗於巴比倫。\\\hline
 (47章)&主前609年，埃及進攻非利士(王下23:29)。\\\hline
 (48章)&摩押的傾覆。\\\hline
 (49:1-6)&亞捫被擄。\\\hline
 (49:7-22)&以東的敗亡(參俄巴底亞書)。\\\hline
 (49:23-27)&亞蘭國的災禍。\\\hline
 (49:28-33)&主前599年，巴比倫攻打亞拉伯的夏瑣，基達。\\\hline
 (49:34-39)&以攔的結局。\\\hline 
(50-51章)&主前539年，巴比倫亡于瑪代波斯，遙指末後大災難世界的毀滅。\\\hline
\end{tabular}
\end{table}
 

\subsection{神賜的異象和比喻}
\begin{table}[H]
\centering
\begin{tabular}{p{3.5cm}|p{2.50cm}|p{11.0cm}}\hline\hline
\multicolumn{1}{c|}{異象}&\multicolumn{1}{c|}{经文}&\multicolumn{1}{c}{寓意}\\\hline
杏樹枝& 1:11-12(民17)&復活的能力，留意保守我的話使得成就\\\hline 
從北而傾燒開的鍋&1:13-16&災禍從北方的巴比倫而來，侵犯，掠奪。\\\hline
曠野熱風&4:11-31&仇敵如雲,地震山搖,天也無光,刀劍害及性命\\\hline
不服法術的毒蛇&8:13-17&從但那裡敵人的馬噴氣發嘶聲，全地震動。他們來吞滅這地和居民\\\hline
甩機弦帳篷毀壞&10:17-22&大擾亂從北方出來使猶大城邑變為荒涼，成為野狗的住處\\\hline
拉出將宰的羊&12:1-4&將惡人拉出來,叫他們等候殺戮的日子\\\hline
與步行的人,馬賽跑&12:5-6&你弟兄和你父家都用奸詐待你，雖向你說好話，你也不要信他們。\\\hline
林中獅子斑點鷙鳥&12:7-9&發聲攻擊神\\\hline
麻帶糜爛&13:1-11&以色列全家和猶大全家也必像這腰帶變為無用\\\hline
酒滿壇&13:12-14&使一切居民都酩酊大醉，彼此相碰，不憐憫，以致滅絕他們\\\hline
旱災&14:2-5&猶大悲哀，城門衰敗沒有水，因為無雨降在地上，地都乾裂。\\\hline
不可娶妻，生兒養女&16:1-4&像糞土不得葬埋，被刀劍和饑荒滅絕。屍首給飛鳥和野獸做食物。\\\hline
窯匠摶泥&18:1-4&在神手中順服，接受神的工作，不順服，必受虧損，另作用途。\\\hline 
毀窯匠的瓦瓶&19&陀斐特和欣嫩子谷，反倒稱為殺戮谷，耶和華打碎這民和這城\\\hline 
巴施戶珥易名&20:1-6&耶和華必使你自覺驚嚇，你也必使眾朋友驚嚇\\\hline
牧人趕散羊大衛苗裔&20:1-6&牧人趕散之羊必從北方和各國旋返,大衛苗裔掌王權\\\hline
\multirow{2}{5.8cm}{二筐無花果}&\multirow{2}{2.8cm}{24:1-7}&極好初熟的無花果是順服神的兒女(林後3:16)。\\\cline{3-3}
&&極壞的無花果，不順服神的兒女。\\\hline
列國民喝憤怒的酒&25:15-29&命刀劍臨到地上一切的居民,向地上一切的居民呐喊，像踹葡萄的一樣。\\\hline
索與軛&27:1-11&西底家登基,喻列國必侍巴比倫王\\\hline
木軛換鐵軛&28&哈拿尼雅折耶利米項上之軛以示列國得釋,木軛換鐵軛，哈拿尼雅當年必亡\\\hline

十七舍克勒銀子贖地&32:9&是神救贖撒但俘虜的表記。十七是復活之數，一月十七日清晨主復活。\\\hline
法老宮門的磚石&43:8-13&預示巴比倫滅埃及。\\\hline
不倒空的酒&48:11-12&舊生命未經過十字架的破碎，新生命不能出來。\\\hline
\end{tabular}
\end{table}

\subsection{先知11次禱告}
\begin{table}[H]
\centering
\begin{tabular}{p{2.cm}|p{10.50cm}}\hline\hline
\multicolumn{1}{c|}{经文}&\multicolumn{1}{c}{内容}\\\hline
 (1:6)&為傳達神的話語呼求神。\\\hline
 (3:10)&為民敗壞光景哀歎。\\\hline
 (10:6-10)&稱頌神的偉大，眾人當敬畏神。\\\hline
 (10:23-25)&求神從寬，行在祂命定的道路中。\\\hline
 (12:1-4)&求問神惡人為何亨通。\\\hline
 (14:7-22)&求神同住不離開，為百姓認罪哀求神施憐憫。\\\hline
 (16:15-18)&向神訴冤，為忠心傳神的話遭淩辱，不想宴樂，痛苦不止而靜坐。\\\hline
 (17:12-18)&求神醫治背道的病，拯救離棄活泉的惡行。\\\hline
 (18:19-23)&求神管教敵擋的兒女。\\\hline
 (20:7-13)&向神訴明裡外之火的追逼光景。\\\hline
 (32:16-25)&稱頌神的偉大與超越。―― 張向晨《聖經六十六卷》 
\\\hline
\end{tabular}
\end{table}



\subsection{耶利米書预言}
%\begin{changemargin}{-1cm}{-1cm}
\begin{sidewaystable} % <-- HERE
\begin{tikzpicture}[thick,snake=zigzag, line before snake = 5mm, line after snake = 5mm]

     % draw 1st horizontal line 
     \draw (-1.,1.55) node[below=3pt] {} node[above=3pt] {以赛亚};
     %\draw (0.,2) node[below=3pt] {} node[above=3pt] {1};                         
     \draw (0.5,2) node[below=3pt] {} node[above=3pt] {};                         
    % \draw (1.,2) node[below=3pt] {} node[above=3pt] {3};                         
    % \draw (1.5,2) node[below=3pt] {} node[above=3pt] {4};                                        
     \draw (2.,2) node[below=3pt] {} node[above=3pt] {};                              
  %  \draw (2.5,2) node[below=3pt] {} node[above=3pt] {6};                  
  %  \draw (21.5,2) node[below=3pt] {} node[above=3pt] {7};                                                       
 
     \draw (0.5,4.4) node[below=3pt] {} node[above=3pt] {亚};   
    \draw (0.5,4.) node[below=3pt] {} node[above=3pt] {兰};  
     \draw (0.5,3.6) node[below=3pt] {} node[above=3pt] {被};   
    \draw (0.5,3.2) node[below=3pt] {} node[above=3pt] {亚};   
     \draw (0.5,2.8) node[below=3pt] {} node[above=3pt] {述};   
    \draw (0.5,2.4) node[below=3pt] {} node[above=3pt] {灭}; 
   
     \draw (1.5,4.4) node[below=3pt] {} node[above=3pt] {北};   
    \draw (1.5,4.) node[below=3pt] {} node[above=3pt] {国};  
     \draw (1.5,3.6) node[below=3pt] {} node[above=3pt] {被};   
    \draw (1.5,3.2) node[below=3pt] {} node[above=3pt] {亚};   
     \draw (1.5,2.8) node[below=3pt] {} node[above=3pt] {述};   
    \draw (1.5,2.4) node[below=3pt] {} node[above=3pt] {掳}; 
                               
     \draw (2.5,4.8) node[below=3pt] {} node[above=3pt] {南};   
    \draw (2.5,4.4) node[below=3pt] {} node[above=3pt] {国};                                
     \draw (2.5,4.0) node[below=3pt] {} node[above=3pt] {被};   
    \draw (2.5,3.6) node[below=3pt] {} node[above=3pt] {巴};                                
     \draw (2.5,3.2) node[below=3pt] {} node[above=3pt] {比};   
    \draw (2.5,2.8) node[below=3pt] {} node[above=3pt] {伦};                                    
    \draw (2.5,2.4) node[below=3pt] {} node[above=3pt] {掳};                                                   

  % \draw (3.5,3.6) node[below=3pt] {} node[above=3pt] {天使印};                                         
   \draw (3.5,3.2) node[below=3pt] {} node[above=3pt] {南};                                         
   \draw (3.5,2.8) node[below=3pt] {} node[above=3pt] {国};                                            
   \draw (3.5,2.4) node[below=3pt] {} node[above=3pt] {归};                                       
   \draw (3.5,2.0) node[below=3pt] {} node[above=3pt] {回};                                            


   \draw (6,3.6) node[below=3pt] {} node[above=3pt] {弥};                                         
   \draw (6,3.2) node[below=3pt] {} node[above=3pt] {赛};                                         
   \draw (6,2.8) node[below=3pt] {} node[above=3pt] {亚};                                            
   \draw (6,2.4) node[below=3pt] {} node[above=3pt] {一次};                                          
  \draw (6,2.0) node[below=3pt] {} node[above=3pt] {来};                                                
%  \draw (6,1.6) node[below=3pt] {} node[above=3pt] {弥};                                         

   \draw (7,3.2) node[below=3pt] {} node[above=3pt] {耶};                                         
   \draw (7,2.8) node[below=3pt] {} node[above=3pt] {稣};                                         
   \draw (7,2.4) node[below=3pt] {} node[above=3pt] {的};                                            
   \draw (7,2.) node[below=3pt] {} node[above=3pt] {工作};                                          
 % \draw (7,1.6) node[below=3pt] {} node[above=3pt] {死};                                                

   \draw (8,3.2) node[below=3pt] {} node[above=3pt] {耶};                                         
   \draw (8,2.8) node[below=3pt] {} node[above=3pt] {稣};                                         
   \draw (8,2.4) node[below=3pt] {} node[above=3pt] {被};                                            
   \draw (8,2.) node[below=3pt] {} node[above=3pt] {钉死};                                          
 
  \draw (9,3.2) node[below=3pt] {} node[above=3pt] {耶};                                         
   \draw (9,2.8) node[below=3pt] {} node[above=3pt] {稣};                                         
   \draw (9,2.4) node[below=3pt] {} node[above=3pt] {复};                                            
   \draw (9,2.) node[below=3pt] {} node[above=3pt] {活};                                          



   \draw (10,3.2) node[below=3pt] {} node[above=3pt] {圣};                                         
   \draw (10,2.8) node[below=3pt] {} node[above=3pt] {灵};                                         
   \draw (10,2.4) node[below=3pt] {} node[above=3pt] {降临};                                            
   \draw (10,2.) node[below=3pt] {} node[above=3pt] {临};                                          


  
   \draw (12,3.2) node[below=3pt] {} node[above=3pt] {教};                                         
   \draw (12,2.8) node[below=3pt] {} node[above=3pt] {会};                                            
   \draw (12,2.4) node[below=3pt] {} node[above=3pt] {被};                                          
   \draw (12,2.) node[below=3pt] {} node[above=3pt] {提};                                                



   \draw (15,3.6) node[below=3pt] {} node[above=3pt] {弥};                                         
   \draw (15,3.2) node[below=3pt] {} node[above=3pt] {赛};                                         
   \draw (15,2.8) node[below=3pt] {} node[above=3pt] {亚};                                            
   \draw (15,2.4) node[below=3pt] {} node[above=3pt] {二次};                                          
  \draw (15,2.) node[below=3pt] {} node[above=3pt] {来};                                                


   \draw (16,3.6) node[below=3pt] {} node[above=3pt] {以色};                                         
   \draw (16,3.2) node[below=3pt] {} node[above=3pt] {列和};                                         
   \draw (16,2.8) node[below=3pt] {} node[above=3pt] {犹大};                                            
   \draw (16,2.4) node[below=3pt] {} node[above=3pt] {二次};                                          
  \draw (16,2.) node[below=3pt] {} node[above=3pt] {归回};                                                


  \draw (17,3.2) node[below=3pt] {} node[above=3pt] {一};                                         
   \draw (17,2.8) node[below=3pt] {} node[above=3pt] {次};                                            
   \draw (17,2.4) node[below=3pt] {} node[above=3pt] {复};                                          
  \draw (17,2.) node[below=3pt] {} node[above=3pt] {活};                                                

  \draw (18,3.2) node[below=3pt] {} node[above=3pt] {审};                                         
   \draw (18,2.8) node[below=3pt] {} node[above=3pt] {判};                                            
   \draw (18,2.4) node[below=3pt] {} node[above=3pt] {列};                                          
  \draw (18,2.) node[below=3pt] {} node[above=3pt] {国};                                                


  \draw (19.2,3.2) node[below=3pt] {} node[above=3pt] {千};                                                
  \draw (19.2,2.8) node[below=3pt] {} node[above=3pt] {禧};                                                   \draw (19.2,2.4) node[below=3pt] {} node[above=3pt] {年};                                                    \draw (19.2,2.) node[below=3pt] {} node[above=3pt] {开始};                                                

  \draw (21.5,4.8) node[below=3pt] {} node[above=3pt] {撒旦};                                                
  \draw (21.5,4.4) node[below=3pt] {} node[above=3pt] {出};                                                   \draw (21.5,4.) node[below=3pt] {} node[above=3pt] {无底坑};                                                    \draw (21.5,3.6) node[below=3pt] {} node[above=3pt] {诱惑};                                                
\draw (21.5,3.2) node[below=3pt] {} node[above=3pt] {列国};                                               
\draw (21.5,2.8) node[below=3pt] {} node[above=3pt] {围攻};                                                    \draw (21.5,2.4) node[below=3pt] {} node[above=3pt] {耶路};                                                
\draw (21.5,2.) node[below=3pt] {} node[above=3pt] {撒冷};                                               


\draw (22.5,4.8) node[below=3pt] {} node[above=3pt] {所有};                                                
\draw (22.5,4.4) node[below=3pt] {} node[above=3pt] {的人};                                                   \draw (22.5,4.) node[below=3pt] {} node[above=3pt] {复活};                                                   
\draw (22.5,3.6) node[below=3pt] {} node[above=3pt] {白色};                                                
\draw (22.5,3.2) node[below=3pt] {} node[above=3pt] {大};                                                   \draw (22.5,2.8) node[below=3pt] {} node[above=3pt] {宝};                                                    \draw (22.5,2.4) node[below=3pt] {} node[above=3pt] {座};                                                
\draw (22.5,2.0) node[below=3pt] {} node[above=3pt] {审判};                                                

  \draw (23.5,3.2) node[below=3pt] {} node[above=3pt] {进};                                                
  \draw (23.5,2.8) node[below=3pt] {} node[above=3pt] {入};                                                
    \draw (23.5,2.4) node[below=3pt] {} node[above=3pt] {永};                                                          
   \draw (23.5,2.) node[below=3pt] {} node[above=3pt] {世};                                                          

 %
      \draw (10,3.8) -- (10,5); 
      \draw (12,3.8) -- (12,5); 
\draw[>=triangle 45, <->] (10,5) -- (12,5); 
  \draw (10.5,5.0) node[below=3pt] {} node[above=3pt] {教};                                                
  \draw (10.9,5.0) node[below=3pt] {} node[above=3pt] {会};                                                
    \draw (11.3,5.0) node[below=3pt] {} node[above=3pt] {时};                                                          
   \draw (11.7,5.0) node[below=3pt] {} node[above=3pt] {期};                                                          


      \draw (15,4) -- (15,5); 
\draw[>=triangle 45, <->] (12,5) -- (15,5); 
  \draw (13.,5.0) node[below=3pt] {} node[above=3pt] {圣};                                                
  \draw (13.4,5.0) node[below=3pt] {} node[above=3pt] {徒};                                                
    \draw (13.8,5.0) node[below=3pt] {} node[above=3pt] {得};                                                          
   \draw (14.2,5.0) node[below=3pt] {} node[above=3pt] {赏};                                                          
  \draw (13.,4.36) node[below=3pt] {} node[above=3pt] {地};                                                
  \draw (13.4,4.36) node[below=3pt] {} node[above=3pt] {上};                                                
    \draw (13.8,4.36) node[below=3pt] {} node[above=3pt] {灾};                                                          
   \draw (14.2,4.35) node[below=3pt] {} node[above=3pt] {难};                                                          


       \draw (-0.5,2) -- (25,2); 
      % \draw (0.,2.1) -- (0,2);       
        \draw (0.5,2.1) -- (0.5,2);    
      % \draw (1.,2.1) -- (1,2);       
        \draw (1.5,2.1) -- (1.5,2);    
      % \draw (2.,2.1) -- (2,2);       
        \draw (3.5,2.1) -- (3.5,2);   
       \draw (16,2.1) -- (16.,2);     
       \draw (17.,2.1) -- (17.,2); 
     \draw (18.,2.1) -- (18.,2); 
       \draw (19.,2.1) -- (19.,2);           
       \draw (21.5,2.1) -- (21.5,2); 
      \draw (22.5,2.1) -- (22.5,2);  
      \draw (23.5,2.1) -- (23.5,2);  
     \iffalse      
      \draw (23.8,2.1) -- (23.8,2);
     \draw (24.8,2.1) -- (24.8,2);
     \draw (25.8,2.1) -- (25.8,2);    
      \fi                                                                   
      \fill (2.5,2)  circle[radius=3pt];
   \fill (6.,2)  circle[radius=3pt]; 
   \fill (15.,2)  circle[radius=3pt];                                           
     \fill (19.2,2)  circle[radius=3pt]; 
     \fill (23.5,2)  circle[radius=3pt];  
\end{tikzpicture}
\end{sidewaystable} 
%\end{changemargin}


